\chapter{Theoretical Basis for the Candidate Criteria}\label{app:criteria-theory}
% Nguon: 03_phan-tich/tieu-chi/2026-07-05_danh-sach-tieu-chi-chinh-thuc_v5.0_final_user.md
%        (14 tieu chi ung vien tu literature, TRUOC phong van) + 1 tieu chi bo sung
%        (Software Compatibility) sau phong van, nhung o day CHI trinh bay co so
%        LY THUYET/hoc thuat cho tieu chi nay, KHONG dua interpretation cua chuyen gia
%        vao (interpretation cua chuyen gia da nam o Ch4 subsec:interview-outcomes).
%        Viet lai doc lap tu tieng Viet sang tieng Anh, khong dich may.

This appendix sets out the academic and theoretical grounding for the fourteen candidate
criteria that literature synthesis (\Cref{subsec:lit-synthesis}) contributed to the
candidate set, together with one additional criterion, software compatibility, that was
introduced only after the expert interviews of \Cref{subsec:expert-interviews}. The
discussion below is confined to prior research and general consumer-electronics theory; it
deliberately excludes the interview-specific reasoning already reported for the
Vietnamese context in \Cref{subsec:interview-outcomes}, so that the two sources of
justification, literature and fieldwork, remain visibly separate rather than blended into
a single narrative.

\section{Economic Factors and the Budget Constraint}\label{app:theory-economic}

Price occupies a central position in the consumer decision network
\parencite{sonmez_cakir_determination_2020}. \textcite{rau_optimal_2018} add an economic
framing in which price is treated as compensation for brand value, and device value is
understood to depreciate as a function of time. The mechanism behind this depreciation is
the pace of technological turnover: a device purchased at a high price can become
functionally dated within a few years even without any physical fault, simply because the
surrounding software and peripheral ecosystem has moved on. This dynamic explains why, in
terms of relative priority, price behaves as a comparatively strong constraint for the
general consumer segment \parencite{weng_siew_lam_evaluation_2023}. Price-sensitive buyers
are willing to trade away hardware headroom they do not need, for instance accepting a
mid-tier processor generation instead of the newest flagship chip, in exchange for a budget
that still covers the tasks they actually perform \parencite{maghsoudi_hybrid_2026}. Expert
or professional buyers, by contrast, tend to treat price as secondary, accepting a
meaningfully higher cost in return for durability and dependable after-sales support. At
the opposite end of the market, \textcite{jimenez-parra_key_2014} find that among buyers of
remanufactured electronics a low price is the one incentive strong enough to make a buyer
accept the residual risk of using previously owned technology.

\section{Core Hardware Capability}\label{app:theory-hardware}

The processing capacity of the CPU together with the working memory supplied by RAM forms
the microarchitectural basis for a computer's execution capability. Because the CPU
functions as the bottleneck that governs a device's overall responsiveness,
\textcite{kang_study_2022} place it among the core quality attributes, one whose weight in
the purchase decision outweighs that of accompanying after-sales services. Examined from a
behavioural angle, \textcite{liao_determinants_2022} document a shift in how consumers
reason about processors: buyers increasingly favour multi-core architectures chosen not
purely to maximise clock speed but to optimise performance per watt of energy consumed, a
preference aligned with broader sustainability expectations. Alongside the processor, a
large RAM capacity has stopped being an optional upgrade and has instead become a
near-mandatory baseline, a shift driven by the requirement to keep several cloud-based
platforms and resource-intensive applications running at once. This same shift toward
modern hardware also explains why solid-state storage has displaced the traditional hard
disk drive: because an SSD contains no spinning platter or other moving mechanical part, it
removes the risk of data loss from physical shock while travelling and reduces power draw
at the same time.

The graphics processing unit (GPU) exposes a clearer split between customer segments.
Users with specialised needs, such as gaming or design work, require strong graphics
performance and willingly accept the bulkier, noisier cooling system that a discrete GPU
demands. The sentiment analysis of social-media posts conducted by
\textcite{maghsoudi_hybrid_2026}, however, shows that heat output and fan noise are among
the strongest sources of complaint precisely within the office-worker segment. This
rejection follows directly from a physical and an economic trade-off. Physically, a
discrete GPU requires a larger cooling assembly and draws more power, which increases the
device's thickness and weight and shortens battery life, all of which work against the
mobility that office use depends on. Economically, because office tasks such as word
processing and spreadsheet work cannot make use of a discrete GPU's capacity, fitting one
becomes a case of over-specification: it raises the budget constraint discussed in
\Cref{app:theory-economic} without returning a proportional benefit
\parencite{rau_optimal_2018}. These two mechanisms together explain why a component that
represents, on paper, superior technical capability is nonetheless declined by the general
buyer segment in favour of integrated graphics.

\section{Portability and Durability}\label{app:theory-portability}

The ability to operate away from a fixed power source makes battery life one of the
attributes buyers weigh most heavily when choosing a laptop.
\textcite{weng_siew_lam_evaluation_2023} report that, among technical specifications,
battery life carries the single highest weight in their model. Buyers oriented toward a
more sustainable lifestyle place a comparable emphasis on efficient batteries as a way of
lengthening the replacement cycle and limiting electronic waste
\parencite{liao_determinants_2022}. Manufacturers nonetheless face a physical trade-off in
delivering this: battery capacity scales with the overall weight of the device.

On design and weight, the evidence runs against the assumption that younger consumers chase
a slim form factor regardless of cost. The data reported by
\textcite{weng_siew_lam_evaluation_2023} show practicality dominating aesthetics: their
student respondents place durability, battery life, and processing power ahead of other
attributes, pushing a slim or fashion-oriented design to the bottom of the priority
ranking. \textcite{zakeri_decision_2023} reinforce this reading by treating durability as
the property that determines how far a device's usable lifecycle can be extended. This
body of evidence explains why buyers accept a somewhat heavier machine, or pay a premium
for an impact-resistant aluminium-alloy shell, in preference to a thin plastic shell that is
light but comparatively fragile \parencite{liao_determinants_2022}.

\section{Interface and Interaction Experience}\label{app:theory-interface}

Across long working sessions, the display and the keyboard are the two attributes that
determine physical comfort and, by extension, how long a session can be sustained without
fatigue. \textcite{maghsoudi_hybrid_2026} record that experienced users hold a demanding
standard for typing feel, so that a keyboard lacking adequate key travel or tactile
feedback draws negative feedback quickly and consistently.

For the display, panel technology matters alongside size. \textcite{liao_determinants_2022}
trace the decline of LCD panels to the arrival of LED backlighting, which can cut power
consumption by as much as 60 percent, a gain that aligns with the broader move toward more
sustainable consumption. Anti-glare coatings and blue-light filtering are likewise becoming
a de facto standard, valued for reducing digital eye strain over extended use. Physical
connectivity ports follow the opposite trajectory: rather than continuing to be upgraded,
they are increasingly folded into a secondary, supplementary feature group
\parencite{sonmez_cakir_determination_2020}, a pattern that reflects the broader shift
toward wireless connectivity. The growth of cloud storage together with a maturing
ecosystem of Bluetooth peripherals has allowed bulkier physical ports, such as optical
drives or wired network sockets, to be removed from many devices altogether, freeing up
space for a more minimalist chassis.

\section{Trust and Emotional Value}\label{app:theory-trust}

Brand is a factor that shapes purchase intention directly, and
\textcite{sonmez_cakir_determination_2020} treat it as a filter through which a buyer
absorbs the physical and technical properties of a product without evaluating each one
individually. Read through a behavioural-economics lens, this filtering role is a response
to information asymmetry: because the general buyer lacks the technical background to
assess a chip, a cooling assembly, or the durability of a motherboard directly, brand
functions as a signalling mechanism that stands in for that missing expertise. Apple is
the clearest illustration of the mechanism: office buyers rarely investigate transistor
counts or the microarchitecture of an M-series chip, and instead purchase on the strength
of a logo that is trusted to guarantee smooth performance, display quality, and battery
life. A less established brand offering comparable specifications on paper typically
struggles to justify the same price, precisely because it lacks this trust filter to
convert raw technical numbers into a sense of purchase safety. Buyers who are strongly
attached to a given brand tend, as a consequence, to be comparatively insensitive to price
movements elsewhere in the market \parencite{rau_optimal_2018}. This trust function is not
always carried by the original manufacturer's logo: in the market for used or
remanufactured devices, it is instead the reputation of the distributor or refurbisher that
carries this role, countering the buyer's default assumption that a used device is more
likely to fail \parencite{jimenez-parra_key_2014}.

After-sales service is treated by buyers as a form of insurance against risk.
\textcite{sostar_assessment_2023} find that a flexible warranty policy reduces the fear of
an expensive component failure and, in doing so, builds a longer-lasting emotional
attachment between the buyer and the product. Although after-sales support is often
classified as a secondary attribute relative to hardware, \textcite{maghsoudi_hybrid_2026}
show that a service failure can trigger a rapid wave of negative sentiment on social media,
so that a single after-sales incident is capable of outweighing the goodwill built by
genuine hardware improvements.

\section{Software Compatibility (Added After the Expert Interviews)}\label{app:theory-software}

Unlike the fourteen criteria discussed above, software compatibility was not part of the
literature-derived candidate set; it entered the candidate set only after the expert
interviews of \Cref{subsec:expert-interviews}, for reasons specific to the Vietnamese office
context that are reported in \Cref{subsec:interview-outcomes}. Independently of that
fieldwork, the wider consumer-electronics literature also treats platform or operating
system as a measurable purchase attribute in its own right, rather than a background
assumption. \textcite{guan_cognitive_2022} build a laptop-recommendation model in which
operating system, coded across Linux, OS X, Windows 7, and Windows 8, is included as a
standalone attribute, and derive distinct, internally consistent preference values for each
option, treating platform choice as directly comparable to attributes such as CPU or RAM
rather than as an incidental detail. \textcite{weng_siew_lam_evaluation_2023}, summarising a
smartphone-selection study by Chakraborty et al., similarly report operating system as the
second most heavily weighted criterion overall, behind brand and ahead of user experience,
functionality, affordability, and design. Taken together, these two sources provide an
independent theoretical basis for treating software or platform compatibility as a
legitimate candidate criterion for personal computing devices, a basis that stands apart
from, and precedes, the Windows-based accounting and enterprise-software lock-in argument
documented from the Vietnamese expert interviews in \Cref{subsec:interview-outcomes}.

\section{Summary}\label{app:theory-summary}

\Cref{tab:criteria-theory-summary} lists each candidate criterion alongside a short
description and the sources on which the discussion above draws.

\begin{table}[htbp]
  \centering
  \caption{Candidate criteria and their supporting literature.}
  \label{tab:criteria-theory-summary}
  \begin{tabular}{@{}p{2.6cm}p{5.0cm}p{6.2cm}@{}}
    \toprule
    \textbf{Criterion} & \textbf{Description} & \textbf{Sources} \\
    \midrule
    Price & Purchase cost of the device & Sönmez Çakır \& Pekkaya (2020); Rau \& Fang (2018); Jiménez-Parra et al. (2014); Lam et al. (2023); Maghsoudi et al. (2026) \\
    CPU & Processor capability & Kang et al. (2022); Liao \& Chuang (2022); Lam et al. (2023); Maghsoudi et al. (2026); Rau \& Fang (2018) \\
    RAM & Memory capacity & Sönmez Çakır \& Pekkaya (2020); Lam et al. (2023); Maghsoudi et al. (2026); Rau \& Fang (2018) \\
    SSD & Storage capacity and type & Sönmez Çakır \& Pekkaya (2020); Lam et al. (2023); Maghsoudi et al. (2026); Rau \& Fang (2018); Liao \& Chuang (2022) \\
    GPU & Graphics processing capability & Maghsoudi et al. (2026); Rau \& Fang (2018) \\
    Display & Screen size and panel technology & Sönmez Çakır \& Pekkaya (2020); Liao \& Chuang (2022); Maghsoudi et al. (2026); Lam et al. (2023) \\
    Keyboard & Typing feel and build quality & Maghsoudi et al. (2026); Liao \& Chuang (2022) \\
    Connectivity ports & Range of physical ports & Maghsoudi et al. (2026); Sönmez Çakır \& Pekkaya (2020) \\
    Design and shell material & Aesthetics and chassis material & Lam et al. (2023); Kang et al. (2022); Maghsoudi et al. (2026) \\
    Size and weight & Portability & Lam et al. (2023); Zakeri et al. (2023); Maghsoudi et al. (2026) \\
    Battery life & Runtime on battery & Lam et al. (2023); Liao \& Chuang (2022); Maghsoudi et al. (2026); Kang et al. (2022); Rau \& Fang (2018) \\
    Durability & Shock resistance and build life & Zakeri et al. (2023); Lam et al. (2023); Liao \& Chuang (2022) \\
    Brand & Manufacturer reputation & Sönmez Çakır \& Pekkaya (2020); Rau \& Fang (2018); Jiménez-Parra et al. (2014); Lam et al. (2023) \\
    After-sales & Warranty and customer support & Šostar \& Ristanović (2023); Maghsoudi et al. (2026); Kang et al. (2022) \\
    Software compatibility & Operating-system and platform fit \newline \textit{(added after expert interviews)} & Guan \& Yuen (2022); Lam et al. (2023) \\
    \bottomrule
  \end{tabular}
\end{table}
